\documentclass[11pt]{article}

\usepackage[a4paper,margin=1in]{geometry}
\usepackage{amsmath,amssymb,amsthm,mathtools}
\usepackage{bm}
\usepackage{enumitem}
\usepackage{hyperref}

\hypersetup{
  colorlinks=true,
  linkcolor=blue,
  citecolor=blue,
  urlcolor=blue
}

\newcommand{\dd}{\mathrm{d}}
\newcommand{\E}{\mathbb{E}}
\newcommand{\R}{\mathbb{R}}
\newcommand{\1}{\mathbf{1}}
\newcommand{\calU}{\mathcal{U}}
\newcommand{\calS}{\mathcal{S}}
\newcommand{\calX}{\mathcal{X}}
\newcommand{\calN}{\mathcal{N}}
\newcommand{\Pgen}{P_{\mathrm{gen}}}
\newcommand{\Pcond}{P_{\mathrm{cond}}}
\newcommand{\sgn}{\operatorname{sgn}}
\newcommand{\card}[1]{\left|#1\right|}

\title{MCMCMC Set Heat-Bath Sampling for CDet Integrals}
\author{}
\date{\today}

\begin{document}
\maketitle

\tableofcontents

\section{Purpose}

This note describes the ``many-configuration Markov-chain Monte Carlo''
(MCMCMC) / set heat-bath algorithm that is currently being tested as a
replacement for the old order-update histogram MCMC.  The goal is to estimate
several signed CDet integrals at different perturbation orders by running a
single Markov chain on a larger ``envelope'' configuration space.

The main idea is:
\begin{quote}
At every Monte Carlo step we keep an envelope of \(N\) vertices.  Inside this
envelope we can see many lower-order configurations at once, namely all
subsets \(S\subset\{1,\ldots,N\}\).  Instead of proposing one add/remove move
and accepting or rejecting it, we compute the heat-bath flux of all visible
subsets, accumulate their signed contribution exactly, choose one subset with
probability proportional to its absolute flux, and regenerate the missing
vertices.
\end{quote}

This is rejection-free in the usual Metropolis sense: every heat-bath update
produces a new envelope.  The analog of an ``acceptance ratio'' is therefore
identically one.  The useful diagnostics are instead the chosen subset
cardinality, the retained fraction \(\card{S}/N\), and the autocorrelation of
the measured order signals.

\section{Target signed integrals}

Let a single vertex be denoted by
\[
  X\in\calX,
\]
where \(\calX\) is equipped with a normalized base measure \(\mu\).  In the
current C++ code the Monte Carlo coordinate is
\[
  z=(s,r_1,r_2,a)\in [0,1)^4,
\]
with the map
\[
  \tau=\beta s,\qquad
  x=\lfloor N_1 r_1\rfloor,\qquad
  y=\lfloor N_2 r_2\rfloor,\qquad
  \alpha=\lfloor 2a\rfloor \in\{0,1\}.
\]
Thus the sampler may be viewed either as sampling the continuous unit cube, or
as sampling a mixed object consisting of imaginary time, lattice site, and
replica index.

For each perturbation order \(u\) we have a signed real integrand
\[
  C_u(X_1,\ldots,X_u),
\]
and the order-\(u\) integral is
\[
  I_u
  =
  \int_{\calX^u}
  C_u(X_1,\ldots,X_u)\,
  \dd\mu(X_1)\cdots \dd\mu(X_u).
\]
The code path currently evaluates \(C_u\) through
\[
  \texttt{integrand\_temp\_typed}
\]
after building the replica Green-function Chebyshev tables.

The algorithm estimates ratios
\[
  \frac{I_u}{I_r},
\]
where \(r\) is a chosen reference/normalization order.  If \(I_r\) is known
from a separate calculation, all \(I_u\) can be put on an absolute scale.

\section{Envelope configurations and subsets}

Choose an envelope order \(N\).  In the current test for measured orders
\[
  \calU=\{2,3,4\},
\]
we choose
\[
  N=5.
\]
This is important: using \(N>u_{\max}\) guarantees that after selecting any
measured subset there is at least one missing vertex to regenerate, so the
chain can actually move.  If one includes the full set \(u=N\), one must add
an additional refresh mechanism or explicitly exclude the full-set choice, as
done in the reference FastFeynmanDiagrammatics code.

An envelope configuration is
\[
  V=(X_1,\ldots,X_N)\in\calX^N.
\]
For a label subset \(S\subseteq [N]\), where \([N]=\{1,\ldots,N\}\), write
\[
  V_S=\{X_i:i\in S\},\qquad
  \bar S=[N]\setminus S,\qquad
  u=\card{S}.
\]
The visible order-\(u\) configurations inside \(V\) are all subsets with
\(\card{S}=u\).  The binomial number
\[
  B_{N,u}=\binom{N}{u}
\]
will appear repeatedly.  It compensates for the fact that one envelope
contains \(B_{N,u}\) labeled order-\(u\) subsets.

\section{The conditional generator}

\subsection{Most general form}

For every subset \(S\) that may be kept, the algorithm needs a normalized
conditional density
\[
  Q_{\bar S|S}(V_{\bar S}\mid V_S)
\]
for generating the missing vertices \(V_{\bar S}\) from the retained vertices
\(V_S\).  The only strict requirements are:
\begin{align}
  Q_{\bar S|S}(V_{\bar S}\mid V_S) &\ge 0,\\
  \int_{\calX^{N-\card{S}}}
  Q_{\bar S|S}(V_{\bar S}\mid V_S)\,
  \dd\mu(V_{\bar S}) &= 1,
\end{align}
and the density must be evaluable for the old complement as well as samplable
for the new complement.

This is the first major freedom of the algorithm.  Any normalized conditional
generator gives an unbiased ratio estimator, provided the same generator is
used consistently in the flux and in the envelope update.

\subsection{Pair-seed construction used by FastFeynmanDiagrammatics}

The reference implementation constructs \(Q_{\bar S|S}\) from a normalized
pair seed
\[
  q(X'\mid X),\qquad
  \int_{\calX}q(X'\mid X)\,\dd\mu(X')=1.
\]
Starting from a set \(B\), a new vertex is generated by first choosing a
parent vertex in \(B\) and then drawing from the pair seed.  With uniform
parent selection,
\[
  q(X\mid B)
  =
  \frac{1}{\card{B}}\sum_{Y\in B} q(X\mid Y).
\]

For two disjoint sets \(A\) and \(B\), define
\[
  \Pgen(A\mid B)
\]
as the probability density of generating all vertices in \(A\), in any order,
starting from \(B\).  The recursion is
\begin{align}
  \Pgen(\varnothing\mid B)&=1,\\
  \Pgen(A\mid B)
  &=
  \frac{1}{\card{A}\card{B}}
  \sum_{X\in A}\sum_{Y\in B}
  q(X\mid Y)\,
  \Pgen(A\setminus\{X\}\mid B\cup\{X\}).
  \label{eq:pgen-recursion}
\end{align}
The factor \(1/\card{A}\) averages over which vertex of \(A\) was generated
first; the factor \(1/\card{B}\) is the uniform parent choice.

For a fixed envelope \(V\), it is convenient to write
\[
  \Pcond^V(A)
  :=
  \Pgen(V_A\mid V_{\bar A}).
\]
Then for \(A\ne\varnothing\),
\begin{align}
  \Pcond^V(A)
  =
  \frac{1}{\card{A}\card{\bar A}}
  \sum_{i\in A}\sum_{j\in \bar A}
  q(X_i\mid X_j)\,
  \Pcond^V(A\setminus\{i\}).
  \label{eq:pcond-recursion}
\end{align}
This dynamic program costs \(O(N^2 2^N)\), rather than summing explicitly over
all possible generation orders.

In the FastFeynmanDiagrammatics code this is exactly the role of the files
\[
  \texttt{compute\_seeds.hpp},\qquad
  \texttt{compute\_conditional\_probabilities.hpp}.
\]
The seed matrix is
\[
  \texttt{seeds}_{ij}=q(X_i\mid X_j),
\]
and \(\Pcond^V(A)\) is computed recursively for all subsets \(A\).

\subsection{Generalized parent choice}

Uniform parent choice is not required.  One can replace it by weights
\[
  \pi(Y\mid B),\qquad
  \pi(Y\mid B)\ge 0,\qquad
  \sum_{Y\in B}\pi(Y\mid B)=1,
\]
so that
\[
  q(X\mid B)=\sum_{Y\in B}\pi(Y\mid B)q(X\mid Y).
\]
Then the recursion becomes
\[
  \Pgen(A\mid B)
  =
  \frac{1}{\card{A}}
  \sum_{X\in A}
  q(X\mid B)\,
  \Pgen(A\setminus\{X\}\mid B\cup\{X\}).
\]
This may be useful if some vertices are more natural anchors than others.

\subsection{Uniform-seed specialization}

The simplest choice is the flat seed
\[
  q(X'\mid X)=1
\]
with respect to the normalized base measure \(\mu\).  Then
\[
  Q_{\bar S|S}(V_{\bar S}\mid V_S)=1,
  \qquad
  \Pcond^V(\bar S)=1.
\]
This is the specialization currently implemented in
\[
  \texttt{include/mcmcmc\_sampler.h}.
\]
It is useful as a correctness baseline because it should agree with plain
uniform Monte Carlo.  It is not expected to be the optimal production
proposal.

\section{Fluxes and the local heat bath}

Let \(\lambda_u>0\) be an order-dependent guiding parameter.  For a measured
subset \(S\) with \(u=\card{S}\in\calU\), define the signed flux
\[
  \Phi_S(V)
  =
  \lambda_u\,
  C_u(V_S)\,
  Q_{\bar S|S}(V_{\bar S}\mid V_S).
  \label{eq:signed-flux}
\]
For the pair-seed construction,
\[
  Q_{\bar S|S}(V_{\bar S}\mid V_S)
  =
  \Pcond^V(\bar S).
\]

The local absolute heat-bath normalization is
\[
  W(V)
  =
  \sum_{\substack{S\subseteq [N]\\ \card{S}\in\calU}}
  \frac{|\Phi_S(V)|}{\binom{N}{\card{S}}}.
  \label{eq:local-weight}
\]
The factor \(1/\binom{N}{\card{S}}\) is the same ``one binomial coefficient''
factor used in the reference implementation.

At a fixed envelope \(V\), the heat bath chooses a subset with probability
\[
  p(S\mid V)
  =
  \frac{|\Phi_S(V)|/\binom{N}{\card{S}}}{W(V)}.
  \label{eq:subset-choice}
\]
After selecting \(S_\star\), the next envelope is produced by retaining
\(V_{S_\star}\) and sampling a new complement from
\[
  Q_{\bar S_\star|S_\star}(\cdot\mid V_{S_\star}).
\]

The signed virtual time associated with a subset is
\[
  t_S(V)
  =
  \frac{\Phi_S(V)}{W(V)}.
  \label{eq:times}
\]
The algorithm accumulates all \(t_S(V)\), not only the chosen subset.  This
is the ``many-configuration'' part: one Markov step produces exact signed
contributions from all visible subsets.

\section{Detailed balance and rejection-free updates}

The Markov chain on envelopes has stationary density proportional to \(W(V)\)
with respect to \(\mu^N\):
\[
  \pi(V)=\frac{W(V)}{Z_W},\qquad
  Z_W=\int_{\calX^N} W(V)\,\dd\mu^N(V).
\]
To see detailed balance, fix a retained subset \(S\).  Suppose \(V\) and
\(V'\) share the same retained vertices \(V_S=V'_S\), while their complements
are different.  The probability flux through this retained subset is
proportional to
\[
  \pi(V)\,
  p(S\mid V)\,
  Q_{\bar S|S}(V'_{\bar S}\mid V_S).
\]
Using Eq.~\eqref{eq:subset-choice}, this is proportional to
\[
  W(V)\,
  \frac{|\lambda_u C_u(V_S)Q_{\bar S|S}(V_{\bar S}\mid V_S)|}
       {\binom{N}{u}W(V)}
  Q_{\bar S|S}(V'_{\bar S}\mid V_S).
\]
Since \(V_S=V'_S\), this equals
\[
  \frac{|\lambda_u C_u(V_S)|}{\binom{N}{u}}\,
  Q_{\bar S|S}(V_{\bar S}\mid V_S)\,
  Q_{\bar S|S}(V'_{\bar S}\mid V_S),
\]
which is symmetric under \(V\leftrightarrow V'\).  Thus the heat-bath move is
already balanced; there is no Metropolis rejection step.

Therefore the Metropolis-style acceptance ratio is
\[
  A_{\mathrm{heatbath}}=1.
\]
Useful diagnostics are instead:
\begin{itemize}[leftmargin=2em]
  \item chosen order histogram \(N_u=\#\{t:\card{S_t}=u\}\),
  \item retained fraction \(\card{S_t}/N\),
  \item final \(\lambda_u\) values after warmup,
  \item autocorrelation of the de-\(\lambda\) order signals.
\end{itemize}

\section{Estimators}

Define the de-\(\lambda\) order signal
\[
  A_u(V)
  =
  \frac{1}{\binom{N}{u}}
  \sum_{\substack{S\subseteq[N]\\\card{S}=u}}
  \frac{t_S(V)}{\lambda_u}.
  \label{eq:order-signal}
\]
Using Eqs.~\eqref{eq:signed-flux} and \eqref{eq:times},
\[
  W(V)A_u(V)
  =
  \frac{1}{\binom{N}{u}}
  \sum_{\card{S}=u}
  C_u(V_S)
  Q_{\bar S|S}(V_{\bar S}\mid V_S).
\]
Integrating over the envelope gives
\begin{align}
  \int W(V)A_u(V)\,\dd\mu^N(V)
  &=
  \frac{1}{\binom{N}{u}}
  \sum_{\card{S}=u}
  \int C_u(V_S)
  \left[
  \int Q_{\bar S|S}(V_{\bar S}\mid V_S)\,\dd\mu(V_{\bar S})
  \right]
  \dd\mu(V_S)
  \nonumber\\
  &=
  \frac{1}{\binom{N}{u}}
  \sum_{\card{S}=u}
  I_u
  =
  I_u.
\end{align}
Therefore
\[
  \E_{\pi}[A_u]=\frac{I_u}{Z_W}.
\]
The unknown normalization \(Z_W\) is common to all orders, so ratios are
\[
  \frac{I_u}{I_r}
  =
  \frac{\E_\pi[A_u]}{\E_\pi[A_r]}.
\]

Equivalently, if one accumulates the raw signal
\[
  M_u(V)
  =
  \frac{1}{\binom{N}{u}}
  \sum_{\card{S}=u}t_S(V),
\]
then
\[
  \E_\pi[M_u]=\frac{\lambda_u I_u}{Z_W},
\]
and the ratio is
\[
  \frac{I_u}{I_r}
  =
  \frac{\E_\pi[M_u]/\lambda_u}{\E_\pi[M_r]/\lambda_r}.
\]
The current C++ code stores the de-\(\lambda\) signal \(A_u\).

\section{Order normalization}

Choose a reference order \(r\in\calU\).  The Monte Carlo directly gives
\[
  R_{u/r}
  =
  \frac{I_u}{I_r}.
\]
If \(I_r\) is known analytically or from a separate high-precision calculation,
then
\[
  I_u=R_{u/r}I_r.
\]
In the current test,
\[
  r=2.
\]
This matches the idea of using a low order as a normalization sector.

\section{The \texorpdfstring{$\lambda_u$}{lambda} parameters}

The \(\lambda_u\)'s are positive guiding parameters.  They do not change the
answer if divided out in the estimator.  They only change how frequently each
order is selected and therefore affect variance and autocorrelation.

Current implementation:
\begin{enumerate}[leftmargin=2em]
  \item Initialize from
  \[
    \texttt{MCMCMC\_LAMBDAS}.
  \]
  \item During warmup only, count the chosen orders in blocks.
  \item Compare the measured chosen-order fraction \(\hat p_u\) with a target
  \(p_u^\star\).  The current target is uniform over the measured orders:
  \[
    p_u^\star=\frac{1}{|\calU|}.
  \]
  \item Update
  \[
    \lambda_u
    \leftarrow
    \lambda_u
    \exp\left[
      \eta\log\left(\frac{p_u^\star}{\hat p_u}\right)
    \right],
  \]
  where \(\eta=\texttt{MCMCMC\_LAMBDA\_ADAPT\_RATE}\).
  \item Freeze the \(\lambda_u\)'s for production.  Warmup samples are
  discarded.
\end{enumerate}

Other valid choices:
\begin{itemize}[leftmargin=2em]
  \item target one special normalization sector more often,
  \item target the highest order more often,
  \item use a Robbins--Monro decreasing adaptation rate,
  \item precompute \(\lambda_u\) in a separate pilot run and keep them fixed,
  \item use the median \(\lambda_u\) across many processes, as in the
  FastFeynmanDiagrammatics production examples.
\end{itemize}

\section{Guiding functions versus physical observables}

The direct version used in the current code sets the flux function equal to
the signed physical CDet integrand:
\[
  \Phi_S(V)=\lambda_u C_u(V_S)Q_{\bar S|S}(V_{\bar S}\mid V_S).
\]
This is the simplest choice and is ideal for comparing to plain Monte Carlo.

A more general production version can use a positive guiding function
\[
  G_u(V_S)\approx |C_u(V_S)|
\]
inside the absolute heat-bath flux, while measuring the physical observable by
reweighting.  This is conceptually what the FastFeynmanDiagrammatics examples
do with quantities such as \(\Sigma_a(S)\): the heat bath is driven by a
positive proxy, and the desired signed/complex observable is accumulated as a
ratio against that proxy.

This is another major freedom.  The guiding function should be cheap enough
to evaluate and close enough to the magnitude of the target observable to
reduce variance.

\section{Proposal design for the present CDet vertex}

The current uniform proposal is
\[
  q(z'\mid z)=1,\qquad z,z'\in[0,1)^4.
\]
It is correct but not optimal.  A more physical proposal should factorize as
\[
  q(z'\mid z)
  =
  q_\tau(s'\mid s)\,
  q_{\mathrm{space}}(r'_1,r'_2\mid r_1,r_2)\,
  q_{\mathrm{rep}}(a'\mid a).
\]

\subsection{Time proposal}

Since Green functions are strongly structured in time difference, use a
periodic local proposal on \([0,\beta)\).  Examples:
\[
  q_\tau(\tau'\mid\tau)
  \propto
  \frac{1}{1+d_\beta(\tau',\tau)/\tau_0},
\]
or a Cauchy-like kernel
\[
  q_\tau(\tau'\mid\tau)
  \propto
  \frac{1}{1+\left(d_\beta(\tau',\tau)/\tau_0\right)^2},
\]
where \(d_\beta\) is the periodic distance on the imaginary-time circle.
The FastFeynmanDiagrammatics finite-\(\beta\) examples use a logarithmic time
proposal of this spirit, controlled by a scale parameter called
\(\texttt{time\_scale}\).

\subsection{Space proposal}

For lattice sites use a torus distance \(d_L\) and a decaying probability
mass
\[
  q_{\mathrm{lat}}(R'\mid R)
  =
  \frac{w(d_L(R',R))}{\sum_{R''}w(d_L(R'',R))}.
\]
A useful family is
\[
  w(d)=
  \frac{1}{1+c_2(d/\ell)^2+c_4(d/\ell)^4}.
\]
This is close to the form used by the FastFeynmanDiagrammatics
\(\texttt{lattice\_proposer}\), where the scale is controlled by
\(\texttt{space\_scale}\).

If the sampler still works in unit-cube variables \(r_1,r_2\), one can first
choose a lattice site \(R'=(x',y')\), then draw
\[
  r'_1\sim \mathrm{Unif}\left[\frac{x'}{N_1},\frac{x'+1}{N_1}\right),
  \qquad
  r'_2\sim \mathrm{Unif}\left[\frac{y'}{N_2},\frac{y'+1}{N_2}\right).
\]
The density with respect to \(\dd r'_1\dd r'_2\) is then
\[
  N_1N_2\,q_{\mathrm{lat}}(R'\mid R).
\]

\subsection{Replica proposal}

For two replicas, a simple kernel is
\[
  q_{\mathrm{rep}}(\alpha'\mid\alpha)
  =
  (1-p_{\mathrm{flip}})\1_{\alpha'=\alpha}
  +
  p_{\mathrm{flip}}\1_{\alpha'\ne\alpha}.
\]
Uniform replica proposal corresponds to \(p_{\mathrm{flip}}=1/2\).  A local
proposal with smaller \(p_{\mathrm{flip}}\) may be better if same-replica
correlations dominate.

In unit-cube variables, after choosing \(\alpha'\in\{0,1\}\), draw
\[
  a'\sim \mathrm{Unif}\left[\frac{\alpha'}{2},\frac{\alpha'+1}{2}\right).
\]
The density with respect to \(\dd a'\) is
\[
  2\,q_{\mathrm{rep}}(\alpha'\mid\alpha).
\]

\subsection{Mixtures for ergodicity}

To guarantee broad support and avoid overly local trapping, use a mixture
\[
  q_{\mathrm{mix}}(X'\mid X)
  =
  (1-\epsilon)q_{\mathrm{local}}(X'\mid X)
  +
  \epsilon q_{\mathrm{uniform}}(X'),
\]
with \(0<\epsilon\ll 1\).  The density used in \(\Pcond\) must be the same
mixture density used to generate new vertices.

\section{Algorithm summary}

Given a current envelope \(V=(X_1,\ldots,X_N)\):

\begin{enumerate}[leftmargin=2em]
  \item For all relevant subsets \(S\subseteq[N]\), compute
  \(C_{\card{S}}(V_S)\).
  \item Compute \(Q_{\bar S|S}(V_{\bar S}\mid V_S)\).  In the pair-seed
  version, build the seed matrix and compute all \(\Pcond^V(A)\) recursively.
  In the current uniform version this factor is one.
  \item Form signed fluxes
  \[
    \Phi_S(V)=\lambda_{\card{S}}C_{\card{S}}(V_S)Q_{\bar S|S}(V_{\bar S}\mid V_S).
  \]
  \item Form the local normalization
  \[
    W(V)=\sum_S \frac{|\Phi_S(V)|}{\binom{N}{\card{S}}}.
  \]
  \item Accumulate all signed times
  \[
    t_S(V)=\Phi_S(V)/W(V).
  \]
  \item Add the de-\(\lambda\) order signal
  \[
    A_u(V)=\binom{N}{u}^{-1}\sum_{\card{S}=u}t_S(V)/\lambda_u.
  \]
  \item Choose \(S_\star\) with probability
  \[
    \frac{|\Phi_{S_\star}(V)|/\binom{N}{\card{S_\star}}}{W(V)}.
  \]
  \item Keep \(V_{S_\star}\) and regenerate \(V_{\bar S_\star}\) from
  \(Q_{\bar S_\star|S_\star}\).
\end{enumerate}

\section{Mapping to the current C++ implementation}

The implementation supports both uniform and learned factorized-histogram
seeds.  Only the uniform-seed special case has
\[
  Q_{\bar S|S}=1.
\]
For a histogram seed the conditional regeneration density is evaluated
explicitly and retained in the flux.  A positive uniform mixture supplies
full support; production sampling freezes the learned histogram.
For speed, the code does not call the CDet integrand independently for every
subset.  At one envelope \(V\), it builds the full \(N\times N\) up/down
matrices once, computes all principal minors once, and runs an all-subset
z-transform
\[
  \{c(S):S\subseteq[N]\}.
\]
The subset integrand is then read as
\[
  C_{\card{S}}(V_S)
  =
  \frac{(-2)^{\card{S}}}{\card{S}!\,N_1N_2}\,c(S).
\]
This is the CDet analog of the FastFeynmanDiagrammatics idea that exponentially
many biased configurations can be obtained at little extra cost.
Important code locations:
\begin{itemize}[leftmargin=2em]
  \item \(\texttt{include/cdet\_integrand\_evaluator.h}\): wraps the CDet
  integrand evaluator.
  \item \(\texttt{include/mcmcmc\_sampler.h}\): implements the envelope
  heat-bath, the \(\lambda_u\) warmup, chain parallelism, and the plain MC
  benchmark.
  \item \(\texttt{src/main.cpp}\): loads parameters, builds Green-function
  tables, runs MCMCMC, and prints ratios.
  \item \(\texttt{params.txt}\): controls \(N\), measured order range,
  \(\lambda_u\), warmup length, production length, and benchmark sample count.
\end{itemize}

\subsection{The \(\texttt{load\_data}\) stage}

The name \(\texttt{load\_data}\) is slightly historical.  In this branch it
does not read precomputed Green functions from disk.  It constructs the
resident replica Green-function Chebyshev tables for the single parameter set
specified by \(\texttt{params.txt}\), then stores those tables globally for all
later CDet and MCMCMC evaluations.

There are two ways this stage can be reached:
\begin{enumerate}[leftmargin=2em]
  \item \(\texttt{LOAD\_DATA\_WORKFLOW=resident}\), the default path used by
  local runs.  The process builds every table at startup, optionally writes or
  reads the lightweight calculator cache, and keeps the final Chebyshev
  coefficients resident in memory.
  \item \(\texttt{LOAD\_DATA\_WORKFLOW=load\_chunks}\), the HPC path.  Earlier
  array jobs have already sampled the Chebyshev nodes, transformed site-pair
  blocks to coefficients, written corner values, and produced a
  \(\texttt{READY}\) marker.  The main run checks that the manifest matches the
  current parameters, then either loads the coefficient chunks into resident
  tables or opens them through the disk-backed LRU table backend.
\end{enumerate}
Both paths produce the same logical objects.  The only difference is whether
the expensive node sampling and coefficient transform happen in this process or
were split into chunked jobs before the MCMCMC run starts.

The inputs that define the tables are
\[
  N_1,\; N_2,\; \beta,\; \mu,\; \alpha_{\uparrow},\; \alpha_{\downarrow},\;
  \texttt{PBC},\; A,\; \texttt{CHEB\_NU},\; \texttt{CHEB\_NV}.
\]
Here the entanglement subregion \(A\) is selected by
\(\texttt{SUBREGION\_MASK}\).  For \(\texttt{origin\_square}\), the radius is
\(\texttt{SUBREGION\_SQUARE\_RADIUS}\); for all other masks the radius is
ignored.  The spin-dependent one-body chemical potentials used in the free
Green functions are
\[
  \mu_{\uparrow}=\mu-\alpha_{\uparrow},\qquad
  \mu_{\downarrow}=\mu-\alpha_{\downarrow}.
\]
For each table kind the code builds a \(\texttt{ReplicaGreenCalculator}\).  This
calculator constructs the square-lattice one-body Hamiltonian
\[
  h_{ij}=-\mu_\sigma\delta_{ij}
  -t\,\delta_{\langle i,j\rangle},
  \qquad t=1,
\]
with the requested boundary condition, diagonalizes it as \(h=V\Lambda V^T\),
builds the subregion mask, and forms the two sector mixing matrices
\[
  S_m=V^T\Omega_m^{-1}V,
  \qquad m=0,1.
\]
For the ordinary \(Z_1\) subtraction tables the mask is forced to
\(A=\varnothing\), so \(S_m=\mathbf 1\); for the twisted tables the mask is the
chosen entanglement subregion.  If
\(\texttt{LOAD\_DATA\_CALCULATOR\_CACHE}\neq\texttt{off}\), this diagonalized
calculator state may be reused, but the cache is checked against
\(N_1,N_2,\beta,\mu_\sigma,\texttt{PBC}\), the mask, the radius, and the real
type before it is accepted.

There are two precision choices, and they play different roles:
\begin{itemize}[leftmargin=2em]
  \item \(\texttt{LOAD\_DATA\_REAL}\) chooses the real type used while building
  the Chebyshev samples and solving the stable one-particle Green-function
  systems.  Supported values are \(\texttt{MP}\), \(\texttt{double}\),
  \(\texttt{long\_double}\), and \(\texttt{double\_double}\).
  \item \(\texttt{VECTOR\_REAL}\) chooses the resident table type used during
  Monte Carlo lookup.  Supported values are \(\texttt{double}\),
  \(\texttt{long\_double}\), and \(\texttt{double\_double}\).
\end{itemize}
Thus the code may build a table in one precision and store it in another.  The
conversion happens immediately after table construction.  If
\(\texttt{GREEN\_TABLE\_DROP\_BITS}>0\), the resident mantissas are truncated
after conversion; this is a controlled table-compression/noise experiment, not
a change of the underlying theory.  The determinant and z-transform arithmetic
then use \(\texttt{INTEGRAND\_REAL}\), which may differ from both construction
and resident-table precision.

The stable complex linear solves can also monitor the conditioning of the
scaled middle matrices.  The control
\(\texttt{LOAD\_DATA\_CONDITION\_CHECK}\) is \(\texttt{auto}\) by default:
the solver estimates the reciprocal one-norm condition number \(r_{\rm cond}\)
and writes a bounded number of warnings to \(\texttt{stderr}\) only when
\[
  r_{\rm cond}
  \le
  \texttt{LOAD\_DATA\_CONDITION\_RCOND\_WARN}.
\]
The default threshold is \(10^{-12}\).  Setting the mode to \(\texttt{off}\)
removes the check, while \(\texttt{strict}\) turns the same event into an
exception.  For the Green-function kernel solves the right hand side is
diagonal, so the code recovers the one-norm of \(M^{-1}\) from the columns
already solved for \(M^{-1}D_{\rm rhs}\); this adds only \(O(N_{\rm site}^2)\)
work on top of the \(O(N_{\rm site}^3)\) solve.  In a LAPACK double-precision
solve with a non-diagonal right hand side, the fallback estimator is LAPACK's
\(\texttt{zgecon}\).

For each spin, \(\texttt{load\_data}\) builds two kinds of tables:
\begin{enumerate}[leftmargin=2em]
  \item Twisted \(n=2\) replica tables with the chosen subregion \(A\):
  \[
    \texttt{g\_table\_up},\qquad \texttt{g\_table\_down}.
  \]
  These are the tables for the unnormalized twisted object
  \(\widetilde Z_2[A]\).
  \item Ordinary thermal tables with \(A=\varnothing\):
  \[
    \texttt{g\_table\_ordinary\_up},\qquad
    \texttt{g\_table\_ordinary\_down}.
  \]
  These supply the \(Z_1\) subtraction.  The evaluator calls them with
  replica count \(1\), and the normalized \(n=2\) integrand subtracts twice
  the ordinary contribution.
\end{enumerate}
Thus the full resident set is
\[
  \texttt{twisted\_up},\quad
  \texttt{twisted\_down},\quad
  \texttt{ordinary\_up},\quad
  \texttt{ordinary\_down}.
\]
When \(\alpha_\uparrow=\alpha_\downarrow\), the up and down Hamiltonians are
identical, and the implementation reuses the same table object for the two
spins.  This is only a storage optimization; it does not change the spin labels
used later in the determinant.

Each resident table is a \(\texttt{PiecewiseChebyshevTable}\).  It stores four
arrays:
\[
  \texttt{upper\_coeffs},\quad
  \texttt{lower\_coeffs},\quad
  \texttt{upper\_corners},\quad
  \texttt{lower\_corners}.
\]
The spatial layout is the full absolute-coordinate table over all replica
indices and all site pairs \((R_1,R_2)\).  No translation-reduced table is used
at this stage.

The coefficient construction is:
\begin{enumerate}[leftmargin=2em]
  \item Build the tensor-product Chebyshev nodes
  \(u_i=\cos((2i+1)\pi/(2N_u))\) and
  \(v_j=\cos((2j+1)\pi/(2N_v))\), where
  \(N_u=\texttt{CHEB\_NU}\) and \(N_v=\texttt{CHEB\_NV}\).
  \item Map every \((u_i,v_j)\in[-1,1]^2\) to one of the two time triangles.
  \item Evaluate the stable sector kernels
  \(\bar G_m(\tau,\tau')\) for \(m=0,1\) on that branch.  The stable formulas
  are the ones described in
  \(\texttt{stable\_replica\_green\_kernel\_note.tex}\); they solve the scaled
  linear systems without explicitly forming large
  \(\exp(\pm\beta\Lambda)\) factors.
  \item Transform the two sector matrices back to replica indices:
  \[
    \bar G_{\alpha\alpha}
      =\frac{\bar G_0+\bar G_1}{2},\qquad
    \bar G_{01}
      =\frac{i(\bar G_0-\bar G_1)}{2},\qquad
    \bar G_{10}
      =-\frac{i(\bar G_0-\bar G_1)}{2}.
  \]
  \item Store one sampled value for each absolute replica/site pair, then apply
  the two-dimensional Chebyshev cosine transform in the node indices.  The
  result is stored in either \(\texttt{upper\_coeffs}\) or
  \(\texttt{lower\_coeffs}\).
\end{enumerate}
For the ordinary \(Z_1\) tables, only the physical replica-\((0,0)\) block is
needed.  The implementation therefore fills the \(N_1N_2\times N_1N_2\)
ordinary block and leaves the rest of the full replica-layout array unused.

The two coefficient arrays are the Chebyshev representations of the two time
triangles:
\[
  \text{upper branch: }\tau>\tau',
  \qquad
  \text{lower branch: }\tau<\tau'.
\]
The square Chebyshev coordinates are mapped to the triangles by
\[
  \tau=\beta\frac{1+u}{2},\qquad
  \tau'=\tau\frac{1+v}{2}
  \quad\text{on the upper branch},
\]
and
\[
  \tau'=\beta\frac{1+u}{2},\qquad
  \tau=\tau'\frac{1+v}{2}
  \quad\text{on the lower branch}.
\]
This is exactly the same split as in the replica Green-function theory:
the upper table represents
\[
  e^{-\tau h}A_m e^{\tau' h},
\]
while the lower table represents
\[
  e^{-\tau h}(A_m-\mathbf 1)e^{\tau' h}.
\]
The corner arrays hold the limiting values at the origin of each triangle.
During lookup, \(\texttt{BRANCH\_EPSILON}\) decides when a requested point is
close enough to that origin to use the stored corner value instead of forming a
nearly singular Chebyshev coordinate.

The lookup map is the inverse of the construction map.  If
\(\tau_1>\tau_2\), the evaluator uses the upper table and sets
\[
  u=2\tau_1/\beta-1,\qquad v=2\tau_2/\tau_1-1.
\]
Otherwise it uses the lower table and sets
\[
  u=2\tau_2/\beta-1,\qquad v=2\tau_1/\tau_2-1.
\]
This includes the equal-time case: the CDet matrix builder selects
\[
  \texttt{use\_upper\_branch}
  =
  (i\neq j)\ \mathrm{and}\ (\tau_i>\tau_j).
\]
Therefore diagonal entries and exact time ties are lower-branch lookups.
This is the correct convention for the equal-time object
\(\bar G(i,\tau;i,\tau^+)\), because the second time is infinitesimally later:
\[
  \bar G^{(m)}(\tau,\tau^+)
  =
  e^{-\tau h}(A_m-\mathbf 1)e^{\tau h}.
\]
The upper branch would instead represent
\(\bar G^{(m)}(\tau^+,\tau)\).  In the code this lower-side limit is encoded by
passing the equal numerical coordinates
\(\tau_1=\tau_2=\tau\) together with
\(\texttt{use\_upper\_branch=false}\); it is not an average over the jump at
equal time.

After these four tables have been built, the integrand layer requires the
resident type to match \(\texttt{VECTOR\_REAL}\).  A twisted evaluation reads
\(\texttt{g\_table\_up/down}\); an ordinary evaluation reads
\(\texttt{g\_table\_ordinary\_up/down}\).  The density actually used by the
normalized \(n=2\) calculation is therefore
\[
  C_{\mathrm{norm}}(S)
  =
  C_{\mathrm{twisted}}^{(n=2)}(S)
  -
  2\,C_{\mathrm{ordinary}}^{(n=1)}(S),
\]
for every subset \(S\) seen by the envelope z-transform.

For the default test
\[
  \calU=\{2,3,4\},\qquad N=5,\qquad r=2,
\]
the code computes all subsets of the five-vertex envelope whose cardinality is
2, 3, or 4.  It prints:
\begin{itemize}[leftmargin=2em]
  \item final \(\lambda_u\)'s for each chain,
  \item chosen order counts,
  \item heat-bath acceptance, equal to one,
  \item mean retained fraction,
  \item de-\(\lambda\) means and ratios relative to the normalization order,
  \item a plain uniform MC benchmark for comparison.
\end{itemize}

\section{Main free choices}

The following choices are algorithmic knobs:
\begin{enumerate}[leftmargin=2em]
  \item Envelope order \(N\).
  \item Measured order set \(\calU\).
  \item Reference order \(r\).
  \item Conditional generator \(Q_{\bar S|S}\), or equivalently pair seed
  \(q(X'\mid X)\) plus parent selection rule.
  \item Whether to use the physical signed integrand directly as the flux, or
  a positive guiding function plus reweighting.
  \item The order-guiding parameters \(\lambda_u\).
  \item Warmup target fractions \(p_u^\star\) and adaptation rate.
  \item Mixture fraction between local and uniform proposal.
  \item Chain count and internal integrand thread count.
  \item Numerical precision for load-data and integrand evaluation.
\end{enumerate}

Only item 4 and item 5 change the variance structure in a deep way.  The
other knobs mainly control balancing, autocorrelation, and computational cost.

\end{document}
